% ============================================================
% Structure, links, and small custom macros
% - Chapter / TOC formatting
% - Notes / comments helpers
% - Colors and hyperlinks
% ============================================================

% Colors (load before `hyperref`)
% `svgnames` gives named colors like `DarkGreen`.
\usepackage[svgnames]{xcolor}
\begin{comment}
  \setcounter{tocdepth}{2}   % 目录深度 (只显示到 section)
\usepackage{titlesec}
\setcounter{secnumdepth}{2}
\usepackage{tocloft} \renewcommand{\cftsecpresnum}{\S} 
\setlength{\cftchapnumwidth}{1.5em}
\setlength{\cftsecnumwidth}{1.5em}
\setlength{\cftsubsecnumwidth}{1.5em}

% Chapter 样式
\titleformat{\chapter}[display]
  {\normalfont\LARGE\bfseries\raggedleft}
  {第\,\thechapter\,章}
  {0pt}
  {
    \vspace*{-1em}
    \begin{tikzpicture}
      \draw[thick] (0,0) -- (\linewidth,0);
    \end{tikzpicture}
    \vspace{1em}
    \huge
  }[\vspace{1em}]
\titlespacing*{\chapter}{0pt}{120pt}{20pt}

\renewcommand\thesection{\arabic{section}}

% Section 样式
\titleformat{\section}[block]
  {\normalfont\large\bfseries\raggedright\sffamily}
  {\S\thesection}{1em}{}
\titlespacing*{\section}{0pt}{1.5em}{1.5em}
\end{comment}

   %TOC depth: show up to `\subsection`
\setcounter{tocdepth}{3}   % 目录深度 (只显示到 section)
% Chapter titles
\usepackage[explicit]{titlesec}
\usepackage{tocloft}
\titleformat{\chapter}[block]
  {\centering\huge\bfseries\vspace{10ex}}
  {第\,\thechapter\,章}
  {10pt}
  {
    #1
  }
  


% Chinese names for standard titles
\renewcommand{\contentsname}{目录}
\renewcommand{\partname}{部分}

% Notes
\usepackage{comment}
\usepackage{multicol}
\newcommand{\Note}[1]{%
  \footnote{\footnotesize #1}
}

% Hyperlinks (keep black by default; switch colors in `preamble/citecolor.tex`)
\usepackage{hyperref}
\hypersetup{
  colorlinks=true,
  linkcolor=black,
  citecolor=black,
  urlcolor=black,
  pdfborder={0 0 0}
}


\usepackage{amsthm}

% 在导言区定义
\newtheoremstyle{postulate}  % 自定义样式名
  {3ex}{3ex}               % 上下间距
  {}{}                       % 正文字体、标签字体
  {\large\bfseries}                % 标题字体
  {.}{ }                     % 标题后标点、标题后空格
  {}                         % 标题格式（空=默认）

\theoremstyle{postulate}
\newtheorem{postulate}{假设}  % 自动编号：假设 1, 假设 2, ...

% 也可以定义其他类型
\newtheorem{axiom}{公理}
\newtheorem{principle}{原理}

\newcommand{\topictitle}[1]{%
  \bigskip
  \begin{center}
    \textbf{\large #1}
  \end{center}
  \medskip
}